\chapter{Tipos de neuronas según su función eléctrica}
\chaptermark{Neuronas como instrumentos}
\label{sec:eetypes}
\begin{chaptergoals}
\item qué instrumentos son las neuronas: integradores, filtros, convertidores, osciladores y cerrojos;
\item cómo una célula se vuelve un resonador pasa banda, o un cerrojo que \nt{SERO} activa;
\item cómo la realimentación hace de neuronas con fugas un integrador sin fugas, que exige ajuste fino;
\item por qué son modos de una célula y no tipos, fijados por los canales iónicos y los de difusión.
\end{chaptergoals}

En el capítulo~\ref{sec:neurontypes} clasificamos las neuronas según el neurotransmisor que libera su salida. Un electrónico las clasificaría de otro modo: según lo que la neurona hace con su señal, es decir, según el instrumento que es. El instrumento principal del libro lo conocemos desde el capítulo~\ref{sec:glutcomputer}: la neurona integradora con fugas~\eqref{eq:glutneuron}, un circuito RC con comparador. Pero en el cerebro hay otros. La tabla~\ref{tab:eetypes} los reúne en cuatro grupos, y casi todos han aparecido ya en el libro.

\begin{center}
\footnotesize
\setlength{\tabcolsep}{3pt}
\renewcommand{\arraystretch}{1.15}
\begin{longtable}{>{\raggedright}p{3.1cm}>{\raggedright}p{3.3cm}>{\raggedright}p{4.7cm}>{\raggedright\arraybackslash}p{2.4cm}}
\caption{Neuronas como instrumentos: nombre, análogo electrónico, qué hace el instrumento y dónde aparece en el libro.}\label{tab:eetypes}\\
\toprule
Neurona & Instrumento & Qué hace & Dónde en el libro \\
\midrule
\endfirsthead
\toprule
Neurona & Instrumento & Qué hace & Dónde en el libro \\
\midrule
\endhead
\bottomrule
\endfoot
\multicolumn{4}{l}{\emph{Filtros e integradores}} \\
neurona integradora con fugas & integrador RC con comparador & suma las entradas en una ventana $\tau$ y dispara en el umbral & cap.~\ref{sec:glutcomputer}, \eqref{eq:glutneuron} \\
detector de coincidencias & AND con ventana temporal & dispara solo si las entradas llegan casi juntas; $\tau$ muy pequeña & caps.~\ref{sec:glutcomputer}, \ref{sec:code}, \eqref{eq:window} \\
neurona fásica & filtro pasa altos, detector de flancos & responde al cambio de la entrada y se calla & cap.~\ref{sec:sensors} \\
neurona adaptativa & filtro pasa altos con control automático de ganancia & la frecuencia cae con una entrada constante & cap.~\ref{sec:sensors}, \eqref{eq:nakarushton} \\
resonador & circuito resonante, filtro pasa banda & responde mejor a una entrada de su propia frecuencia & sección~\ref{sec:resonator} \\
integrador sin fugas & integrador con amplificador operacional & convierte velocidad en posición y la mantiene & sección~\ref{sec:perfectint} \\
\midrule
\multicolumn{4}{l}{\emph{Convertidores}} \\
neurona tónica & convertidor voltaje---frecuencia & la frecuencia sigue linealmente a la entrada y casi no se fatiga & cap.~\ref{sec:code} \\
neurona graduada (sin espigas) & amplificador analógico, búfer & transmite un voltaje continuo a corta distancia & cap.~\ref{sec:sensors} \\
neurona rectificadora & rectificador de media onda & la frecuencia nunca es negativa, $[u]_+$ & caps.~\ref{sec:glutcomputer}, \ref{sec:gaba} \\
par “encendido---apagado” & par de rectificadores, código de dos cables & una responde a “más”, la otra a “menos” & cap.~\ref{sec:glutcomputer}, \eqref{eq:dualrail} \\
neurona con inhibición de derivación & divisor analógico, control de ganancia & divide la entrada en vez de restarle & cap.~\ref{sec:gaba}, \eqref{eq:shunt} \\
\midrule
\multicolumn{4}{l}{\emph{Generadores y tiempo}} \\
marcapasos & oscilador de relajación, multivibrador & emite espigas por sí solo a frecuencia constante & caps.~\ref{sec:serocomputer}, \ref{sec:dofa} \\
marcapasos con entrada & oscilador controlado por voltaje & ritmo propio que la entrada acelera o frena & cap.~\ref{sec:chemoutput} \\
neurona de ráfagas & generador de ráfagas con compuerta & emite ráfagas de espigas rápidas: las “rayas” & cap.~\ref{sec:code}, \eqref{eq:burst} \\
neurona ligada a la fase & detector de fase, lazo de enganche de fase & dispara en una fase determinada de la onda de entrada & caps.~\ref{sec:sensors}, \ref{sec:chemoutput} \\
axón con retardo & línea de retardo & retrasa la señal un tiempo dado & cap.~\ref{sec:glutcomputer}, fig.~\ref{fig:itd} \\
\midrule
\multicolumn{4}{l}{\emph{Memoria y conmutación}} \\
neurona biestable & disparador de Schmitt, cerrojo (latch) & una entrada breve la lleva a un estado duradero & sección~\ref{sec:bistable} \\
neurona con ramas dendríticas & multiplexor, pequeña red de varias capas & una rama deja pasar la entrada solo si hay una entrada de habilitación & cap.~\ref{sec:synthesis} \\
\end{longtable}
\end{center}

Una salvedad importa más que toda la tabla: no son células distintas, sino \emph{modos} distintos. Una misma neurona puede ser un integrador con una entrada lenta y un detector de coincidencias cuando la inhibición ha acortado su $\tau$ (capítulo~\ref{sec:code}); un marcapasos en la vigilia y un receptor silencioso durante el sueño. Qué instrumento resulta lo deciden los canales iónicos de la membrana y los canales de difusión que los ajustan.

\section{Cuatro instrumentos que aún no habían aparecido en el libro}

\subsection{El resonador}
\label{sec:resonator}

Algunas neuronas tienen una corriente lenta que se opone a cualquier cambio de voltaje: si el voltaje sube, la corriente tira de él hacia abajo; si baja, hacia arriba, pero con retardo. Para un electrónico, esa corriente se comporta como una inductancia, y junto con la capacidad de la membrana forma un circuito resonante. La neurona responde con más fuerza a una entrada de cierta frecuencia, normalmente de unos pocos hercios a una o dos decenas, y más débilmente a entradas más lentas y más rápidas~\cite{HutcheonYarom2000}. Es un filtro pasa banda en una sola célula: de una entrada ruidosa selecciona el ritmo de su frecuencia, por ejemplo el ritmo local del capítulo~\ref{sec:gaba}.

\subsection{El integrador sin fugas}
\label{sec:perfectint}

Un integrador con fugas recuerda la entrada solo durante un tiempo $\tau$, decenas de milisegundos. Pero el ojo, para quedarse quieto, necesita recordar su posición durante segundos: la orden de “girar” llega como una velocidad breve, y hay que mantener el ojo en la nueva posición. La red lo resuelve con la misma realimentación que la memoria del capítulo~\ref{sec:glutcomputer}. Si un grupo de neuronas con constante de tiempo $\tau$ se excita a sí mismo con peso $w$, entonces $\tau\,\dd r/\dd t=-r+w\,r+I$, y la constante de tiempo del grupo es
\begin{empheq}[box=\keybox]{equation}
\tau_{\text{ef}} \;=\; \frac{\tau}{1-w} .
\label{eq:perfectint}
\end{empheq}
Con $\tau=0.1\,\text{s}$ y $w=0.995$ se obtienen $20\,\text{s}$: un integrador casi ideal hecho de neuronas que, cada una por sí sola, olvidan en una décima de segundo~\cite{Seung1996}. El precio es el ajuste fino: con $w$ algo mayor que uno el integrador se satura, y con $w$ algo menor olvida enseguida. Por eso este integrador necesita un reajuste constante mediante aprendizaje, como el descrito en el capítulo~\ref{sec:motorlearning}.

\subsection{La neurona biestable}
\label{sec:bistable}

En los capítulos~\ref{sec:glutcomputer} y~\ref{sec:gaba} el biestable se construía con un grupo de neuronas. Pero también hay biestables de una sola célula. Algunas neuronas tienen una corriente que, una vez activada, mantiene por sí misma la excitación: una entrada breve lleva a la neurona a una actividad duradera, una \emph{meseta}, y la inhibición la devuelve al reposo. Es un disparador de Schmitt: la neurona tiene dos estados estables e histéresis entre ellos. Así se comportan, por ejemplo, las neuronas \nt{ACET} que controlan los músculos, y esta propiedad la activa un canal de difusión: la meseta aparece cuando a la neurona le llega serotonina~\cite{Hounsgaard1988}. Una misma célula es un cerrojo con \nt{SERO} y un integrador corriente sin él.

\subsection{Integradores y resonadores}

La teoría divide las células excitables en dos clases~\cite{Izhikevich2007}. Los \emph{integradores} se parecen a un oscilador controlado por voltaje que arranca desde cero: justo por encima del umbral la neurona dispara tan despacio como se quiera, y la frecuencia crece suavemente con la entrada. Una neurona así suma todo lo que llega y transmite bien con su frecuencia la magnitud de la entrada. Los \emph{resonadores} se parecen a un oscilador con frecuencia mínima: no saben disparar despacio; con una entrada en el umbral emiten de inmediato espigas a una frecuencia finita y prefieren una entrada de su propia frecuencia. Los primeros son los instrumentos naturales del código de frecuencia del capítulo~\ref{sec:code}; los segundos, del código temporal.

\begin{keyblock}
\begin{proposition}[Una neurona, muchos instrumentos]
\label{prop:eetypes}
Por su función eléctrica, las neuronas funcionan como integradores con fugas y sin ellas, detectores de coincidencias, filtros pasa altos y pasa banda, convertidores voltaje---frecuencia, rectificadores, divisores, osciladores, líneas de retardo y biestables. Qué instrumento resulta de una célula dada lo fijan sus canales iónicos y los canales de difusión que los ajustan. Los modos en sí están medidos; hasta qué punto usa el cerebro esta reconfiguración para calcular es sobre todo una hipótesis.
\end{proposition}
\end{keyblock}

Para un ingeniero, esto se parece a una matriz lógica programable: el hardware es uno, y el instrumento lo fija la configuración. Solo que aquí la configuración no se cambia al programar el chip, sino sobre la marcha, mediante los canales de difusión lentos de los capítulos~\ref{sec:serocomputer}--\ref{sec:acet}.

\section{Ejercicios}

\begin{exercise}[\lvlA]
\label{ex:18:1}
\emph{Tolerancia del integrador sin fugas.} Neuronas con $\tau=0.1$~s mantienen la posición del ojo según~\eqref{eq:perfectint} durante $T=1$~s con un error de no más del $5\%$ en ambos sentidos. (a)~Halle el intervalo admisible de $w$. (b)~$w$ es la suma de los pesos de $K$ sinapsis, cada una con un error independiente del $10\%$. ¿Con qué $K$ la dispersión de la suma cabe en la tolerancia?
\end{exercise}

\begin{exercise}[\lvlB]
\label{ex:18:2}
\emph{El resonador como circuito.} Describamos la corriente lenta de la sección~\ref{sec:resonator} con una variable $W$: $C\,\dd V/\dd t=-g_LV-g_wW+I$, $\tau_w\,\dd W/\dd t=V-W$. Muestre que es un circuito $RC$ con una rama paralela $R_w=1/g_w$, $L_w=\tau_w/g_w$, y con $C/g_L=10\ms$, $\tau_w=100\ms$, $\alpha=g_w/g_L=4$ halle la frecuencia de resonancia y el cociente $|Z|/|Z(0)|$ en ella.
\end{exercise}

\begin{exercise}[\lvlB]
\label{ex:18:3}
\emph{Cómo empieza a disparar un integrador.} (a)~Neurona $\tau\,\dd V/\dd t=-V+\mu$, $\tau=10\ms$, con umbral $\theta$ y reinicio a cero. ¿Con qué $\mu/\theta-1$ la frecuencia es de $1$~Hz? (b)~¿Qué frecuencia da el modelo $\tau\,\dd v/\dd t=v^2+\mu$ (la espiga es la fuga de $v$ a $+\infty$, con reinicio en $-\infty$) con $\mu=0.01$?
\end{exercise}

\begin{exercise}[\lvlB]
\label{ex:18:4}
\emph{Simulación: integrador y resonador.} El modelo del ejercicio~\ref{ex:18:2} con $g_L=1$, $\tau_m=10\ms$, $\tau_w=100\ms$, umbral $1$ y reinicio $V\to0$ ($W$ no cambia) recibe $\mu=1.5\sin2\pi ft$, $f=1$--$40$~Hz. ¿En qué banda de frecuencias emite espigas la neurona con $\alpha=0$ y con $\alpha=4$? Compare con la condición $1.5\,|Z(f)|>1$.
\end{exercise}

\begin{exercise}[\lvlB]
\label{ex:18:5}
\emph{Un cerrojo de una sola célula.} Neurona con corriente de meseta: $\tau\,\dd r/\dd t=-r+F(wr+I)$, $F(x)=\min(1,\max(0,x))$, $\tau=10\ms$, $w=1.5$, fondo $I=-0.2$. (a)~¿Qué duración mínima debe tener un pulso $I=0.3$ para llevar la neurona de $r=0$ a la meseta? (b)~¿Qué amplitud mínima $B$ debe tener un pulso largo $I=-0.2-B$ para devolverla al reposo?
\end{exercise}

\begin{exercise}[\lvlB]
\label{ex:18:6}
\emph{Autoajuste del integrador.} Al fijar la mirada la entrada es $I=0$, un sensor de deslizamiento da la velocidad de deriva $e=\dd r/\dd t$, y $\dd w/\dd t=-\eta\,e\,r$. Con $w=0.95$, $\tau=0.1$~s, $\langle r^2\rangle=1$, halle el $\eta$ con el que en $60$~s de fijaciones $|1-w|$ entra en la tolerancia del ejercicio~\ref{ex:18:1}. ¿Qué se estropea si también se aprende durante los giros del ojo?
\end{exercise}

\begin{exercise}[\lvlC]
\label{ex:18:7}
\emph{¿Para qué reconfigurar el instrumento?} Un canal de difusión conmuta $\alpha$ del modelo del ejercicio~\ref{ex:18:2} entre $0$ y $4$. ¿Cuántas veces supera el resonador al integrador en relación señal/ruido para una senoide débil en ruido blanco de corriente a $11$~Hz y a $1$~Hz? Invente una tarea en la que convenga precisamente conmutar de modo.
\end{exercise}
